\documentclass[11pt,a4paper]{article}
\usepackage[margin=24mm]{geometry}
\usepackage{amsmath,amssymb,hyperref}
\hypersetup{hidelinks}
\title{Stanley Supplementary Exercise 38\\A reverse-tableau formula for skew tableaux}
\author{Statement: Richard P. Stanley\\Paired original proof: Andrei Okounkov and Grigori Olshanski}
\date{Statement version 4 June 2023; author proof arXiv:q-alg/9605042v1}
\begin{document}
\maketitle
\section*{Original exercise}
Let \(\lambda\vdash n\), \(\mu\vdash k\), and \(\ell=\ell(\lambda)\).
Let \(\operatorname{RT}(\mu,\ell)\) be the set of all reverse SSYT of shape \(\mu\)
and largest part at most \(\ell\). For a square \(u\in\mu\) let \(c(u)\)
denote its content. Write \((n)_k=n(n-1)\cdots(n-k+1)\). Show that
\[
 f^{\lambda/\mu}=\frac{f^\lambda}{(n)_k}
 \sum_{T\in\operatorname{RT}(\mu,\ell)}
 \prod_{u\in\mu}\bigl(\lambda_{T(u)}-c(u)\bigr),
\]
where \(T(u)\) denotes the entry in square \(u\) of \(T\).

\paragraph{Example.} For \(\mu=(2,1)\) we get
\[
 f^{\lambda/21}=\frac{f^\lambda}{(n)_3}
 \left(
 \begin{aligned}
 &\sum_{i<j\le\ell}\lambda_j(\lambda_j-1)(\lambda_i+1)\\
 &+\sum_{i<j\le\ell}\lambda_j(\lambda_i+1)(\lambda_i-1)\\
 &+\sum_{i<j<k\le\ell}\lambda_k(\lambda_j-1)(\lambda_i+1)\\
 &+\sum_{i<j<k\le\ell}\lambda_k(\lambda_i-1)(\lambda_j+1)
 \end{aligned}
 \right).
\]

\section*{Paired author proof}
The official solutions explicitly cite Okounkov and Olshanski,
\emph{Shifted Schur Functions}, Theorems 8.1 and 11.1.
The complete native AMS-TeX source and author PDF are retained with this statement.
The paired source uses the second proof of Theorem 8.1 and the first proof
of Theorem 11.1, including Lemma 11.2 and the required definitions.

\section*{Sources}
Original statement: Stanley, exercise 38, physical page 8,
\url{https://math.mit.edu/~rstan/ec/ch7supp.pdf}.
Official solution pointer: version 24 March 2023, physical page 13,
\url{https://math.mit.edu/~rstan/ec/ch7suppsol.pdf}.
Original author proof:
\url{https://arxiv.org/abs/q-alg/9605042v1}.
\end{document}
